\def\handoutsay{false}\def\handoutfill{false}
\documentclass[12pt]{article}
\newcommand{\handoutrunning}{Lecture 7: Eigenvalues and diagonalization}
\providecommand{\handoutsay}{true}
\usepackage[T1]{fontenc}
\usepackage{lmodern,microtype}
\usepackage[a4paper,margin=15mm,headheight=14pt,headsep=5mm]{geometry}
\usepackage{amsmath,amssymb,amsthm,mathtools}
\usepackage{enumitem,booktabs,needspace}
\usepackage{pifont}
\usepackage[dvipsnames]{xcolor}
\usepackage{fancyhdr}
\usepackage{ifthen}
\usepackage[hidelinks]{hyperref}
\hypersetup{pdfauthor={Chang-Ye Tu}}
\pagestyle{fancy}
\fancyhf{}
\fancyhead[L]{\small Linear Algebra}
\fancyhead[R]{\small\handoutrunning}
\fancyfoot[C]{\small\thepage}
\renewcommand{\headrulewidth}{0.3pt}
\setlength{\parindent}{0pt}
\setlength{\parskip}{4pt plus 1pt minus 1pt}
\setlength{\emergencystretch}{1.5em}
\usepackage{etoolbox}
\newcommand{\handoutdisplayskips}{\setlength{\abovedisplayskip}{7pt plus 2pt minus 2pt}%
  \setlength{\belowdisplayskip}{7pt plus 2pt minus 2pt}}
\apptocmd{\normalsize}{\handoutdisplayskips}{}{\PackageWarning{handout}{could not patch normalsize}}
\handoutdisplayskips
\setlist[enumerate]{label=(\alph*),leftmargin=*,itemsep=2pt,topsep=3pt}
\setlist[itemize]{leftmargin=*,itemsep=2pt,topsep=3pt}
\newcommand{\R}{\mathbb R}
\newcommand{\C}{\mathbb C}
\newcommand{\iu}{\mathrm{i}}
\newcommand{\contradiction}{\mathord{\text{\large\ding{56}}}}
\newcommand{\one}{\mathbf 1}
\DeclareMathOperator{\tr}{tr}
\DeclareMathOperator{\col}{col}
\DeclareMathOperator{\diag}{diag}
\DeclareMathOperator{\rank}{rank}
\DeclareMathOperator\prb{\mathsf{P}}
\DeclareMathOperator\expc{\mathsf{E}}
\DeclareMathOperator\var{var}
\DeclareMathOperator\cov{cov}
\DeclareMathOperator\cor{cor}
\DeclareMathOperator\argmin{argmin}
\DeclareMathOperator\argmax{argmax}
\DeclareMathOperator{\Span}{span}
\theoremstyle{definition}
\newtheorem{theorem}{Theorem}
\newtheorem{proposition}[theorem]{Proposition}
\newtheorem{lemma}[theorem]{Lemma}
\newtheorem{corollary}[theorem]{Corollary}
\newtheorem{definition}[theorem]{Definition}
\newtheorem{example}[theorem]{Example}
\newtheorem{remark}[theorem]{Remark}
\newtheorem{fact}{Fact}
\renewcommand{\thefact}{F\arabic{fact}}
\AddToHook{cmd/subsection/before}{\par\Needspace{9\baselineskip}}
\AddToHook{env/theorem/before}{\par\Needspace{4\baselineskip}}
\AddToHook{env/proposition/before}{\par\Needspace{4\baselineskip}}
\AddToHook{env/lemma/before}{\par\Needspace{4\baselineskip}}
\AddToHook{env/definition/before}{\par\Needspace{4\baselineskip}}
\AddToHook{env/example/before}{\par\Needspace{4\baselineskip}}
\definecolor{saybg}{gray}{0.86}
\NewDocumentEnvironment{say}{+b}
  {\ifthenelse{\equal{\handoutsay}{true}}{\par\Needspace{2\baselineskip}\noindent\colorbox{saybg}{\parbox{\dimexpr\linewidth-2\fboxsep\relax}%
   {\small\sffamily\ding{43}\enspace\textbf{Say.}\enspace#1}}\par}{}}
  {}
\providecommand{\handoutfill}{false}
\newsavebox{\boardbox}
\newlength{\boardht}
\definecolor{boardrule}{gray}{0.55}
\makeatletter
\newcommand{\boardstart}{\par\if@inlabel\leavevmode\par\fi}
\makeatother
\NewDocumentEnvironment{board}{o +b}
  {\boardstart\IfValueT{#1}{\noindent#1\par\nopagebreak}%
   \noindent\begin{lrbox}{\boardbox}\begin{minipage}{\linewidth}#2\end{minipage}\end{lrbox}%
   \setlength{\boardht}{\dimexpr\ht\boardbox+\dp\boardbox\relax}%
   \ifthenelse{\equal{\handoutfill}{true}}%
     {\setlength{\boardht}{1.4\boardht}%
      {\setlength{\fboxsep}{0pt}\setlength{\fboxrule}{0.4pt}\color{boardrule}%
       \framebox[\dimexpr\linewidth-2\fboxrule\relax]{\rule{0pt}{\boardht}}}}%
     {\ifthenelse{\equal{\handoutsay}{true}}%
        {\llap{\color{boardrule}\rule[-\dp\boardbox]{1.5pt}{\boardht}\hspace{4pt}}}{}%
      \usebox{\boardbox}}%
   \par}
  {}
\newcommand{\blank}[1]{\ifthenelse{\equal{\handoutfill}{true}}{\underline{\phantom{#1}}}{#1}}
\newcommand{\lec}[1]{Lecture~#1}
\newcommand{\handouttitle}[2]{%
  \begin{center}
  {\LARGE\sffamily\bfseries #2\par}\vspace{3pt}
  {\sffamily Linear Algebra \textperiodcentered\ \lec{#1}\par}
  \ifthenelse{\equal{\handoutsay}{true}}{\vspace{3pt}{\small\sffamily\colorbox{saybg}{\ding{43}\enspace\textbf{Say.}}\enspace
    teacher's copy: the shaded boxes are not in the students' copy.\par}}{}
  \ifthenelse{\equal{\handoutfill}{true}}{\vspace{3pt}{\small\sffamily Fill-in copy: complete the framed spaces in
    class.\par}}{}
  \end{center}}

\hypersetup{pdftitle={Linear Algebra, Lecture 7: Eigenvalues and diagonalization}}
\newcommand{\boardtop}{\vspace*{-\abovedisplayskip}}
\begin{document}
\handouttitle{7}{Eigenvalues and diagonalization}

\textbf{Last time} (\lec{6}): orthogonality. Results of Lectures~1--5 are cited where they are used; F1--F4:
facts about polynomials, without proof (Chapter~7 of the textbook). \(I=I_n\), \(e_j=\col_j(I)\).

\textbf{Polynomials.} \(\C=\{a+bi:a,b\in\R\}\), \(i^2=-1\); \(zw=0\Rightarrow z=0\) or \(w=0\).
\(f=\sum_kc_kt^k\), \(g=\sum_kd_kt^k\) (\(c_k,d_k\in\C\), \(c_k=d_k=0\) for \(k>m\); \(f\) real if all \(c_k\in\R\)):
\(f=g\) if all \(c_k=d_k\); \(f+g\), \(fg\) have the coefficients \(c_k+d_k\), \(\sum_{i+j=k}c_id_j\); constant term
\(c_0\); \(\deg f=\max\{k:c_k\ne0\}\) (\(f\ne0\)), monic if \(c_{\deg f}=1\). Values \(f(z)=\sum_kc_kz^k\)
(\(z\in\C\)), \(f(A)=\sum_kc_kA^k\) (\(f\) real, \(A^0=I\), \(A^{k+1}=A^kA\)); \(z\) is a root if \(f(z)=0\).
\begin{enumerate}[label=\textbf{F\arabic*.},ref=F\arabic*,leftmargin=*,itemsep=0pt,parsep=0pt,topsep=2pt]
\item\label{f-degree} \(f,g\ne0\Rightarrow fg\ne0\) and \(\deg(fg)=\deg f+\deg g\).
\item\label{f-roots} \(f\ne0\) of degree \(m\) has at most \(m\) roots; so \(f(z)=g(z)\) for infinitely many \(z\)
\(\Rightarrow\) \(f=g\).
\item\label{f-complex-roots} A monic \(f\) of degree \(n\geqslant1\) is \((t-\lambda_1)\cdots(t-\lambda_n)\),
\(\lambda_i\in\C\) unique up to order (fundamental theorem of algebra): roots \(\lambda_1,\dots,\lambda_n\); coefficient of
\(t^{n-1}\): \(-(\lambda_1+\dots+\lambda_n)\); constant term \((-1)^n\lambda_1\cdots\lambda_n\).
\item\label{f-values} \((f+g)(z)=f(z)+g(z)\), \((fg)(z)=f(z)g(z)\) (\(z\in\C\)).
\end{enumerate}

\subsection*{1. Eigenvalues and eigenvectors}

\begin{definition}\label{h-def-eigen}
Let \(A\in\R^{n\times n}\). \(\lambda\in\R\) is an \textbf{eigenvalue} of \(A\) if \(Ax=\lambda x\) for some \(x\in\R^n\),
\(x\ne0\); such \(x\) is an \textbf{eigenvector} belonging to \(\lambda\). The \textbf{eigenspace} of \(\lambda\) is the
subspace \(\mathcal N(A-\lambda I)=\{x\in\R^n:Ax=\lambda x\}\) (\lec{4}, Definition~7: null spaces are subspaces);
its nonzero vectors are the eigenvectors
belonging to \(\lambda\).
\end{definition}
Eigenvalues need not be real numbers: once the characteristic polynomial of \(A\) is defined below, the eigenvalues of
\(A\) are defined as the \(n\) complex roots of that polynomial, listed with multiplicity, and the eigenvalues of this
definition are exactly the real ones among them.

\Needspace{11\baselineskip}
\begin{theorem}\label{h-thm-eig-singular}
Let \(A\in\R^{n\times n}\), \(\lambda\in\R\). Then \(\lambda\) is an eigenvalue of \(A\) \(\iff\)
\(\det(\lambda I-A)=0\) \(\iff\) \(\rank(A-\lambda I)<n\).
\end{theorem}
\begin{proof}
\begin{board}[The chain of equivalences; \(\det(-F)=(-1)^n\det F\) (\lec{5}, Theorem~12(c)):]
\boardtop
\begin{align*}
\lambda\text{ is an eigenvalue}&\iff(A-\lambda I)x=0\text{ for some }x\ne0&&\text{(\lec{1}, Theorem~3(a))}\\
&\iff A-\lambda I\text{ is not invertible}&&\text{(\lec{2}, Theorem~15)}\\
&\iff\rank(A-\lambda I)<n&&\text{(\lec{4}, Theorem~12(d))}\\
&\iff\det(A-\lambda I)=0&&\text{(\lec{5}, Theorem~12(a))}\\
&\iff\det(\lambda I-A)=0&&\text{(\lec{5}, Theorem~12(c))}
\end{align*}
\end{board}
\nopagebreak
\end{proof}

\subsection*{2. The characteristic polynomial}

\textbf{Characteristic polynomial.} For \(A\in\R^{n\times n}\), \(\det(tI-A)\), computed by cofactor expansion
(\lec{5}, Definition~7) with \(t\) as a number, is a real polynomial: the \textbf{characteristic polynomial}
\(p_A(t)=\det(tI-A)\) of \(A\). So \(p_A(c)=\det(cI-A)\) for \(c\in\R\).

\begin{theorem}\label{h-thm-charpoly}
Let \(A=(a_{ij})\in\R^{n\times n}\) and \(\tr A=a_{11}+\dots+a_{nn}\).
\begin{enumerate}
\item \(p_A\) is monic of degree \(n\); its coefficient of \(t^{n-1}\) is \(-\tr A\); its constant term is
\((-1)^n\det A\).
\item \(\lambda\in\R\) is an eigenvalue of \(A\) (Definition~\ref{h-def-eigen}) \(\iff\) \(p_A(\lambda)=0\). So \(A\)
has at most \(n\) such eigenvalues.
\end{enumerate}
\end{theorem}
\begin{proof}
(a)
\begin{itemize}
\item Claim: if \(M\in\R^{m\times m}\) has entries \(\alpha_{ij}+\beta_{ij}t\) (\(\alpha_{ij},\beta_{ij}\in\R\)) and
\(\beta_{ij}\ne0\) occurs in at most \(d\) rows, then \(\det M\) is \(0\) or of degree \(\leqslant d\).
\item Induction on \(m\). \(m=1\): \(\det M=\alpha_{11}+\beta_{11}t\) (\lec{5}, Definition~7: \(\det(a)=a\)).
No \(\beta_{ij}\ne0\) (e.g.\ \(d=0\)): \(\det M\) is a number. Otherwise expand along a row \(i\) with some
\(\beta_{ij}\ne0\) (\lec{5}, Theorem~13(a): expansion along any row); each \(M(i|j)\) has \(\beta\ne0\) in at most \(d-1\) rows, so
\(\det M=\sum_j(-1)^{i+j}(\alpha_{ij}+\beta_{ij}t)\det M(i|j)\) is \(0\) or of degree \(\leqslant d\) (induction,
\ref{f-degree}).
\end{itemize}
Induction on \(n\); \(n=1\): \(p_A=t-a_{11}\) (\lec{5}, Definition~7). For \(n\geqslant2\), as
\((tI-A)(1|1)=tI_{n-1}-A(1|1)\), expanding along row \(1\) (\lec{5}, Theorem~13(a)) gives
\[
p_A=(t-a_{11})\,p_{A(1|1)}+\sum_{j=2}^n(-1)^{1+j}(-a_{1j})\det\bigl((tI-A)(1|j)\bigr)
\]
\begin{itemize}
\item \(j\geqslant2\): row \(j\) of \(tI-A\) loses its only \(t\), so \((tI-A)(1|j)\) has \(t\) in at most \(n-2\) rows;
its determinant is \(0\) or of degree \(\leqslant n-2\) (claim).
\item Induction: \((t-a_{11})\bigl(t^{n-1}-\tr A(1|1)\,t^{n-2}+r\bigr)=t^n-\bigl(a_{11}+\tr A(1|1)\bigr)t^{n-1}+s\), where
\(r\), \(s\) are \(0\) or of degree \(\leqslant n-3\), \(\leqslant n-2\) (\ref{f-degree}); \(a_{11}+\tr A(1|1)=\tr A\).
\item So \(p_A=t^n-(\tr A)\,t^{n-1}+(0\text{ or degree}\leqslant n-2)\).
\end{itemize}
Constant term (every \(n\)): \(p_A(0)=\det(-A)=(-1)^n\det A\) (\lec{5}, Theorem~12(c)).

(b) \(p_A(\lambda)=\det(\lambda I-A)\) and Theorem~\ref{h-thm-eig-singular}; at most \(n\) roots by (a) and
\ref{f-roots}.
\end{proof}

\Needspace{16\baselineskip}
\begin{example}\label{h-exm-charpoly}
Find the eigenvalues and eigenspaces of
\[
A=\begin{pmatrix}1&2&1\\2&0&-2\\-1&2&3\end{pmatrix}
\]
\end{example}
\textit{Solution.}
\begin{board}[\(p_A\), expanding along row \(1\) (\lec{5}, Theorem~13(a)):]
\[
\begin{aligned}
\det(tI-A)&=\det\begin{pmatrix}t-1&-2&-1\\-2&t&2\\1&-2&t-3\end{pmatrix}
=(t-1)\bigl(t(t-3)+4\bigr)+2\bigl(-2(t-3)-2\bigr)-(4-t)\\
&=(t^3-4t^2+7t-4)+(-4t+8)+(t-4)=t^3-4t^2+4t=t(t-2)^2
\end{aligned}
\]
Eigenvalues \(0\), \(2\) (Theorem~\ref{h-thm-charpoly}(b)).
\end{board}
\begin{board}[\(\lambda=0\): reduce \(A-0I=A\); a basis of \(\mathcal N(A)\):]
\[
\begin{gathered}
A\xrightarrow[R_3\leftarrow R_3+R_1]{R_2\leftarrow R_2-2R_1}
\begin{pmatrix}1&2&1\\0&-4&-4\\0&4&4\end{pmatrix}
\xrightarrow{R_3\leftarrow R_3+R_2}
\begin{pmatrix}1&2&1\\0&-4&-4\\0&0&0\end{pmatrix}
\\
\xrightarrow{R_2\leftarrow-\frac14R_2}
\begin{pmatrix}1&2&1\\0&1&1\\0&0&0\end{pmatrix}
\xrightarrow{R_1\leftarrow R_1-2R_2}
\begin{pmatrix}1&0&-1\\0&1&1\\0&0&0\end{pmatrix}
\end{gathered}
\]
Free unknown \(x_3\); special solution \((1,-1,1)'\), a basis of \(\mathcal N(A)\) (\lec{4}, Theorem~8(c)).
\end{board}
\begin{board}[\(\lambda=2\): reduce \(A-2I\); a basis of \(\mathcal N(A-2I)\):]
\[
A-2I=\begin{pmatrix}-1&2&1\\2&-2&-2\\-1&2&1\end{pmatrix}
\xrightarrow[R_3\leftarrow R_3-R_1]{R_2\leftarrow R_2+2R_1}
\begin{pmatrix}-1&2&1\\0&2&0\\0&0&0\end{pmatrix}
\xrightarrow[R_2\leftarrow\frac12R_2]{R_1\leftarrow-R_1}
\begin{pmatrix}1&-2&-1\\0&1&0\\0&0&0\end{pmatrix}
\xrightarrow{R_1\leftarrow R_1+2R_2}
\begin{pmatrix}1&0&-1\\0&1&0\\0&0&0\end{pmatrix}
\]
Special solution \((1,0,1)'\), a basis of \(\mathcal N(A-2I)\) (\lec{4}, Theorem~8(c)).
\end{board}

\begin{definition}\label{h-def-eigenvalues}
Let \(A\in\R^{n\times n}\). By Theorem~\ref{h-thm-charpoly}(a) and \ref{f-complex-roots},
\(p_A=(t-\lambda_1)\cdots(t-\lambda_n)\) with \(\lambda_1,\dots,\lambda_n\in\C\), unique up to order: the
\textbf{eigenvalues} of \(A\), listed with multiplicity; they are the roots of \(p_A\). The \textbf{algebraic
multiplicity} of \(c\in\C\) is \(a_A(c)=\#\{i:\lambda_i=c\}\); \(c\) is \textbf{simple} if \(a_A(c)=1\). The real
\(\lambda_i\) are the eigenvalues of Definition~\ref{h-def-eigen} (Theorem~\ref{h-thm-charpoly}(b)). From now on
``eigenvalue'' has this meaning; eigenvectors and eigenspaces belong to real eigenvalues only.
\end{definition}
In Example~\ref{h-exm-charpoly}: eigenvalues \(0\), \(2\), \(2\); \(a_A(0)=1\), \(a_A(2)=2\).

\Needspace{13\baselineskip}
\begin{example}\label{h-exm-complex}
Find the eigenvalues of
\[
A=\begin{pmatrix}2&-1&0\\1&2&0\\1&1&1\end{pmatrix}
\]
and the eigenspace of each real eigenvalue.
\end{example}
\textit{Solution.} Expanding along column \(3\) (\lec{5}, Theorem~10), and as \((t-2-i)(t-2+i)=(t-2)^2-i^2\),
\[
\det(tI-A)=\det\begin{pmatrix}t-2&1&0\\-1&t-2&0\\-1&-1&t-1\end{pmatrix}=(t-1)\bigl((t-2)^2+1\bigr)
=(t-1)\bigl(t-(2+i)\bigr)\bigl(t-(2-i)\bigr)
\]
Eigenvalues \(1\), \(2+i\), \(2-i\) (Definition~\ref{h-def-eigenvalues}); only \(1\)
is real. For \(\lambda=1\):
\[
A-I=\begin{pmatrix}1&-1&0\\1&1&0\\1&1&0\end{pmatrix}
\xrightarrow[R_3\leftarrow R_3-R_1]{R_2\leftarrow R_2-R_1}
\begin{pmatrix}1&-1&0\\0&2&0\\0&2&0\end{pmatrix}
\xrightarrow{R_3\leftarrow R_3-R_2}
\begin{pmatrix}1&-1&0\\0&2&0\\0&0&0\end{pmatrix}
\]
So \(x_1=x_2=0\), \(x_3\) free: \(\mathcal N(A-I)=\Span\{e_3\}\) (\lec{2}, Theorem~3(b): row operations keep the
solutions); \(2\pm i\) have no eigenvectors in \(\R^3\) (Definition~\ref{h-def-eigenvalues}).

\Needspace{16\baselineskip}
\subsection*{3. Similar matrices}

\begin{theorem}\label{h-thm-similar}
Let \(A,P\in\R^{n\times n}\), \(P\) invertible, and \(B=P^{-1}AP\) (\(B\) is \textbf{similar} to \(A\)). Then
\(p_B=p_A\). So \(A\), \(B\) have the same eigenvalues with the same algebraic multiplicities, \(\tr B=\tr A\) and
\(\det B=\det A\).
\end{theorem}
\begin{proof}
\begin{board}[\(\det(cI-B)=\det(cI-A)\) for \(c\in\R\); \(\det(FG)=\det F\det G\) (\lec{5}, Theorem~12(b)):]
\boardtop
\begin{align*}
\det P^{-1}\det P&=\det(P^{-1}P)=\det I=1&&\text{(\lec{5}, Theorems~12(b), 8(d))}\\
cI-B&=P^{-1}(cI)P-P^{-1}AP=P^{-1}(cI-A)P&&\text{(\lec{1}, Theorem~3(a))}\\
\det(cI-B)&=\det P^{-1}\det(cI-A)\det P=\det(cI-A)&&\text{(\lec{5}, Theorem~12(b))}
\end{align*}
So \(p_B=p_A\) (\ref{f-roots}); eigenvalues, \(\tr\), \(\det\) are read off \(p_A\)
(Definition~\ref{h-def-eigenvalues}, Theorem~\ref{h-thm-charpoly}(a)).
\end{board}
\nopagebreak
\end{proof}

\Needspace{17\baselineskip}
\begin{theorem}\label{h-thm-related}
Let \(A\in\R^{n\times n}\) and \(c\in\R\).
\begin{enumerate}
\item If \(Ax=\lambda x\), \(\lambda\in\R\), \(x\ne0\): \(A^kx=\lambda^kx\) (\(k=0,1,2,\dots\)), \(f(A)x=f(\lambda)x\) for every
real polynomial \(f\), and \((A+cI)x=(\lambda+c)x\).
\item If moreover \(A\) is invertible: \(\lambda\ne0\) and \(A^{-1}x=\lambda^{-1}x\).
\item If \(A\) has the eigenvalues \(\lambda_1,\dots,\lambda_n\), then \(A+cI\) has the eigenvalues
\(\lambda_1+c,\dots,\lambda_n+c\).
\end{enumerate}
\end{theorem}
\begin{proof}
\begin{board}[(a) by induction on \(k\), then \(f(A)x\) and \((A+cI)x\):]
\boardtop
\begin{align*}
A^0x&=x,\qquad A^{k+1}x=A^k(Ax)=\lambda A^kx=\lambda^{k+1}x&&\text{(induction; \lec{1}, Theorem~3(a))}\\
f(A)x&=\textstyle\sum_kc_kA^kx=\sum_kc_k\lambda^kx=f(\lambda)x&&(f=\textstyle\sum_kc_kt^k)\\
(A+cI)x&=Ax+cx=(\lambda+c)x&&\text{(\lec{1}, Theorem~3(a))}
\end{align*}
\end{board}
\begin{board}[(b) \(\lambda\ne0\), then \(A^{-1}x\):]
\boardtop
\begin{align*}
\lambda=0&\implies Ax=0,\ x\ne0\implies A\text{ not invertible }\contradiction&&\text{((b); \lec{2}, Theorem~15)}\\
x&=A^{-1}(Ax)=\lambda A^{-1}x\implies A^{-1}x=\lambda^{-1}x&&\text{(\lec{1}, Theorem~3(a))}
\end{align*}
\end{board}
\begin{board}[(c) For \(s\in\R\), with \(p_A=(t-\lambda_1)\cdots(t-\lambda_n)\) (Definition~\ref{h-def-eigenvalues}):]
\boardtop
\begin{align*}
\det\bigl(sI-(A+cI)\bigr)&=\det\bigl((s-c)I-A\bigr)=p_A(s-c)&&\text{(\lec{1}, Definition~1)}\\
&=\bigl(s-(\lambda_1+c)\bigr)\cdots\bigl(s-(\lambda_n+c)\bigr)&&\text{(\ref{f-values})}
\end{align*}
So \(p_{A+cI}=\bigl(t-(\lambda_1+c)\bigr)\cdots\bigl(t-(\lambda_n+c)\bigr)\) (\ref{f-roots}): the
eigenvalues of \(A+cI\), listed with multiplicity, are \(\lambda_1+c,\dots,\lambda_n+c\)
(Definition~\ref{h-def-eigenvalues}, \ref{f-complex-roots}).
\end{board}
\nopagebreak
\end{proof}

\Needspace{10\baselineskip}
\begin{theorem}\label{h-thm-trace-det}
Let \(A\in\R^{n\times n}\) have the eigenvalues \(\lambda_1,\dots,\lambda_n\). Then \(\tr A=\lambda_1+\dots+\lambda_n\)
and \(\det A=\lambda_1\cdots\lambda_n\). So \(A\) is invertible \(\iff\) \(0\) is not an eigenvalue of \(A\).
\end{theorem}
\begin{proof}
\begin{board}[Coefficients of \(p_A=(t-\lambda_1)\cdots(t-\lambda_n)\) (Definition~\ref{h-def-eigenvalues}):]
\boardtop
\begin{align*}
t^{n-1}&:\quad-\tr A=-(\lambda_1+\dots+\lambda_n)
&&\text{(Theorem~\ref{h-thm-charpoly}(a), \ref{f-complex-roots})}\\
t^0&:\quad(-1)^n\det A=(-1)^n\lambda_1\cdots\lambda_n
&&\text{(Theorem~\ref{h-thm-charpoly}(a), \ref{f-complex-roots})}
\end{align*}
So \(A\) is invertible \(\iff\det A\ne0\) (\lec{5}, Theorem~12(a)) \(\iff\lambda_i\ne0\) for all \(i\)
(\(zw=0\Rightarrow z=0\) or \(w=0\)).
\end{board}
\nopagebreak
\end{proof}
Example~\ref{h-exm-charpoly}: \(\tr A=1+0+3=0+2+2\), \(\det A=0\cdot2\cdot2=0\), so \(A\) is not invertible.
Example~\ref{h-exm-complex}: \(\tr A=2+2+1=1+(2+i)+(2-i)\), \(\det A=1\cdot(2+i)(2-i)=5\).

\subsection*{4. Independence and multiplicities}

\begin{theorem}\label{h-thm-indep}
Let \(v_1,\dots,v_m\) be eigenvectors of \(A\in\R^{n\times n}\) for real eigenvalues \(\lambda_1,\dots,\lambda_m\), with
\(\lambda_i\ne\lambda_j\) for \(i\ne j\). Then \(v_1,\dots,v_m\) are linearly independent. If \(A\) has \(n\) distinct real
eigenvalues, eigenvectors for them form a basis of \(\R^n\).
\end{theorem}
\begin{proof}
Induction on \(m\); \(m=1\): \(c_1v_1=0\), \(v_1\ne0\) \(\Rightarrow\) \(c_1=0\) (\lec{3}, Proposition~2:
\(cv=0\Rightarrow c=0\) or \(v=0\)). Let \(m\geqslant2\) and
\(c_1v_1+\dots+c_mv_m=0\).
\begin{itemize}
\item \(A\) times it minus \(\lambda_m\) times it (\(Av_i=\lambda_iv_i\); \lec{1}, Theorem~3(a)):
\(c_1(\lambda_1-\lambda_m)v_1+\dots+c_{m-1}(\lambda_{m-1}-\lambda_m)v_{m-1}=0\).
\item Induction hypothesis: \(c_i(\lambda_i-\lambda_m)=0\), so \(c_i=0\) for \(i<m\) (\(\lambda_i\ne\lambda_m\)).
\item Then \(c_mv_m=0\) and \(v_m\ne0\), so \(c_m=0\) (\lec{3}, Proposition~2).
\end{itemize}
Last sentence: \(n\) independent vectors in \(\R^n\) form a basis (\lec{4}, Theorem~5(d), as \(\dim\R^n=n\) by
Theorem~4).
\end{proof}

\begin{definition}\label{h-def-multiplicity}
For \(A\in\R^{n\times n}\) and \(c\in\R\), the \textbf{geometric multiplicity} of \(c\) is
\(g_A(c)=\dim\mathcal N(A-cI)=n-\rank(A-cI)\) (\lec{4}, Theorem~11: \(\rank(A-cI)+\dim\mathcal N(A-cI)=n\)). So \(g_A(c)\geqslant1\) \(\iff\) \(c\) is an eigenvalue
(Theorem~\ref{h-thm-eig-singular}).
\end{definition}

\begin{theorem}\label{h-thm-geo-alg}
If \(c\in\R\) is an eigenvalue of \(A\in\R^{n\times n}\), then \(1\leqslant g_A(c)\leqslant a_A(c)\). If \(c\) is
simple, \(g_A(c)=1\).
\end{theorem}
\begin{proof}
\begin{itemize}
\item \(g=g_A(c)\geqslant1\) (Definition~\ref{h-def-multiplicity}). Extend a basis \(x_1,\dots,x_g\) of
\(\mathcal N(A-cI)\) to a basis \(x_1,\dots,x_n\) of \(\R^n\) (\lec{4}, Theorem~5(b): an independent list extends to a
basis; \(\dim\R^n=n\), Theorem~4); \(P=[x_1\ \cdots\ x_n]\) is invertible (\lec{4}, Theorem~12(b),(d): independent
columns \(\iff\) rank \(n\) \(\iff\) invertible).
\item \(Ax_j=cx_j=P(ce_j)\) for \(j\leqslant g\); \(Ax_j=Pm_j\) for some \(m_j\in\R^n\) for \(j>g\), as
\(x_1,\dots,x_n\) span \(\R^n\) (\lec{1}, Theorem~4(b)). So \(AP=PM\) for
\(M=[ce_1\ \cdots\ ce_g\ m_{g+1}\ \cdots\ m_n]\) (\lec{1}, Theorem~4(a)), and \(M=P^{-1}AP\). If \(g=n\): \(M=cI\),
\(p_A=p_M=(t-c)^n\) (Theorem~\ref{h-thm-similar}; \lec{5}, Theorem~9(b): \(\det\) of a diagonal matrix is the product
of its diagonal; \ref{f-values}, \ref{f-roots}), and \(a_A(c)=n\).
\item If \(g<n\), \(M=\left(\begin{smallmatrix}cI_g&K\\0&D\end{smallmatrix}\right)\) with
\(D\in\R^{(n-g)\times(n-g)}\), and for \(s\in\R\), with
\(\det\left(\begin{smallmatrix}F&G\\0&H\end{smallmatrix}\right)=\det F\det H\) (\lec{5}, Theorem~18):
\begin{align*}
sI-M&=\begin{pmatrix}(s-c)I_g&-K\\0&sI_{n-g}-D\end{pmatrix}&&\text{(\lec{5}, Theorem~2(a))}\\
\det(sI-M)&=\det\bigl((s-c)I_g\bigr)\det(sI_{n-g}-D)=(s-c)^gp_D(s)
&&\text{(\lec{5}, Theorems~18, 9(b))}
\end{align*}
So \(p_A=p_M=(t-c)^gp_D\) (Theorem~\ref{h-thm-similar}; \ref{f-values}, \ref{f-roots}).
\item With \(p_D=(t-\delta_1)\cdots(t-\delta_{n-g})\) (Theorem~\ref{h-thm-charpoly}(a), \ref{f-complex-roots}), the
eigenvalues of \(A\) are \(c\) (\(g\) times) and \(\delta_1,\dots,\delta_{n-g}\) (Definition~\ref{h-def-eigenvalues},
\ref{f-complex-roots}); so \(a_A(c)\geqslant g\).
\item \(c\) simple: \(1\leqslant g_A(c)\leqslant a_A(c)=1\).\qedhere
\end{itemize}
\end{proof}
In Example~\ref{h-exm-charpoly}: \(g_A(2)=1<2=a_A(2)\).

\subsection*{5. Diagonalization}

\begin{definition}\label{h-def-diagonalizable}
\(A\in\R^{n\times n}\) is \textbf{diagonalizable} if \(P^{-1}AP=\Lambda\) for an invertible \(P\in\R^{n\times n}\) and a
diagonal \(\Lambda=\diag(\lambda_1,\dots,\lambda_n)\) (diagonal entries \(\lambda_i\in\R\)).
\end{definition}

\begin{theorem}\label{h-thm-diag-criterion}
Let \(A\in\R^{n\times n}\) have the distinct real eigenvalues \(c_1,\dots,c_k\), and \(W_i=\mathcal N(A-c_iI)\).
\begin{enumerate}
\item For an invertible \(P=[p_1\ \cdots\ p_n]\) and \(\Lambda=\diag(\lambda_1,\dots,\lambda_n)\):
\(P^{-1}AP=\Lambda\iff Ap_j=\lambda_jp_j\) (\(j=1,\dots,n\)). If so, \(\lambda_1,\dots,\lambda_n\) are the eigenvalues of
\(A\).
\item The following are equivalent: (i) \(A\) is diagonalizable; (ii) \(\R^n\) has a basis of eigenvectors of \(A\);
(iii) all eigenvalues of \(A\) are real and \(g_A(c_i)=a_A(c_i)\) for each \(i\); (iv)
\(g_A(c_1)+\dots+g_A(c_k)=n\). In (iv), bases of \(W_1,\dots,W_k\) together form a basis of \(\R^n\).
\item If \(A\) has \(n\) distinct real eigenvalues, then \(A\) is diagonalizable.
\end{enumerate}
\end{theorem}
\begin{proof}
(a)
\begin{align*}
P^{-1}AP=\Lambda&\iff AP=P\Lambda&&\text{(\lec{1}, Theorem~3(a))}\\
&\iff\col_j(AP)=\col_j(P\Lambda)\text{ for all }j&&\text{(\lec{1}, Definition~1)}\\
&\iff Ap_j=\lambda_jp_j\text{ for all }j&&\text{(\lec{1}, Theorem~4(a), Proposition~13(a))}\\
P^{-1}AP=\Lambda&\implies p_A=p_\Lambda=(t-\lambda_1)\cdots(t-\lambda_n)
&&\text{(Theorem~\ref{h-thm-similar}; \lec{5}, Theorem~9(b); \ref{f-values}, \ref{f-roots})}
\end{align*}
So \(\lambda_1,\dots,\lambda_n\) are the eigenvalues of \(A\) (Definition~\ref{h-def-eigenvalues}).

(b)
\begin{itemize}
\item (i)\(\Leftrightarrow\)(ii): by (a); \(P\) is invertible \(\iff\) its columns are independent (\lec{4},
Theorem~12(b),(d)) \(\iff\) they form a basis of \(\R^n\) (\lec{4}, Theorem~5(d)), and then they are nonzero
(\lec{3}, Proposition~12(a): a list containing \(0\) is dependent).
\item (i)\(\Rightarrow\)(iii): by (a), the eigenvalues are the real \(\lambda_j\), and
\(a_A(c_i)=d_i=\#\{j:\lambda_j=c_i\}\). The \(j\)th entry of \((\Lambda-c_iI)x\) is \((\lambda_j-c_i)x_j\), so
\(\mathcal N(\Lambda-c_iI)\) has the basis \(\{e_j:\lambda_j=c_i\}\). With \(\rank(P^{-1}FP)=\rank F\) (\lec{4},
Theorem~15(d)) and \(\rank F+\dim\mathcal N(F)=n\) (\lec{4}, Theorem~11), for \(F\in\R^{n\times n}\):
\[
g_A(c_i)=n-\rank(A-c_iI)=n-\rank\bigl(P^{-1}(A-c_iI)P\bigr)=n-\rank(\Lambda-c_iI)=d_i
\]
\item (iii)\(\Rightarrow\)(iv): each eigenvalue is one of \(c_1,\dots,c_k\), so
\(\sum_ig_A(c_i)=\sum_ia_A(c_i)=n\).
\item (iv)\(\Rightarrow\)(ii): let \(w_l\in W_l\), \(w_1+\dots+w_k=0\). For the real polynomial
\(f_i=\prod_{j\ne i}\frac{t-c_j}{c_i-c_j}\) (\(f_1=1\) if \(k=1\)), \(f_i(c_l)=1\) if \(l=i\) and \(0\) if \(l\ne i\)
(\ref{f-values}); \(f(A)w_l=f(c_l)w_l\) (Theorem~\ref{h-thm-related}(a); trivial if \(w_l=0\)). So
\(0=f_i(A)(w_1+\dots+w_k)=f_i(c_1)w_1+\dots+f_i(c_k)w_k=w_i\).
\item So bases \(v_{l1},\dots,v_{lg_l}\) of the \(W_l\) (\(g_l=g_A(c_l)\)) are together independent:
\(\sum_{l,r}\gamma_{lr}v_{lr}=0\) is \(w_1+\dots+w_k=0\) with \(w_l=\sum_r\gamma_{lr}v_{lr}\in W_l\); so all \(w_l=0\),
all \(\gamma_{lr}=0\). These \(n\) eigenvectors form a basis of \(\R^n\) (\lec{4}, Theorem~5(d)).
\end{itemize}
(c) \(n\) eigenvectors for the \(n\) eigenvalues form a basis (Theorem~\ref{h-thm-indep}): (ii) holds.
\end{proof}
Not diagonalizable: Example~\ref{h-exm-charpoly} (\(g_A(2)=1<2=a_A(2)\)) and Example~\ref{h-exm-complex}
(\(2\pm i\) are not real).

\Needspace{21\baselineskip}
\begin{example}\label{h-exm-diagonalize}
Diagonalize \(A\): find an invertible \(P\) and a diagonal \(\Lambda\) with \(P^{-1}AP=\Lambda\), where
\[
A=\begin{pmatrix}5&-6&-6\\-1&4&2\\3&-6&-4\end{pmatrix}
\]
\end{example}
\textit{Solution.}
\begin{board}[\(p_A\): \(\col_2\leftarrow\col_2-\col_3\), then take \(t-2\) out of column \(2\) (\lec{5},
Theorem~8(a),(b), for columns by Theorem~13(b)):]
\[
\det(tI-A)=\det\begin{pmatrix}t-5&6&6\\1&t-4&-2\\-3&6&t+4\end{pmatrix}
=\det\begin{pmatrix}t-5&0&6\\1&t-2&-2\\-3&2-t&t+4\end{pmatrix}
=(t-2)\det\begin{pmatrix}t-5&0&6\\1&1&-2\\-3&-1&t+4\end{pmatrix}
\]
\(R_3\leftarrow R_3+R_2\) (\lec{5}, Theorem~8(b)), expand along column \(2\) (\lec{5}, Theorem~10):
\[
\begin{aligned}
\det(tI-A)&=(t-2)\det\begin{pmatrix}t-5&0&6\\1&1&-2\\-2&0&t+2\end{pmatrix}
=(t-2)\det\begin{pmatrix}t-5&6\\-2&t+2\end{pmatrix}\\
&=(t-2)(t^2-3t+2)=(t-1)(t-2)^2
\end{aligned}
\]
Eigenvalues \(1,2,2\); \(a_A(1)=1\), \(a_A(2)=2\) (Definition~\ref{h-def-eigenvalues}).
\end{board}
\begin{board}[\(\lambda=1\): echelon form of \(A-I\), then back substitution:]
\[
\begin{gathered}
A-I=\begin{pmatrix}4&-6&-6\\-1&3&2\\3&-6&-5\end{pmatrix}
\xrightarrow{R_1\leftrightarrow R_2}
\begin{pmatrix}-1&3&2\\4&-6&-6\\3&-6&-5\end{pmatrix}
\xrightarrow[R_3\leftarrow R_3+3R_1]{R_2\leftarrow R_2+4R_1}
\begin{pmatrix}-1&3&2\\0&6&2\\0&3&1\end{pmatrix}
\\
\xrightarrow{R_2\leftrightarrow R_3}
\begin{pmatrix}-1&3&2\\0&3&1\\0&6&2\end{pmatrix}
\xrightarrow{R_3\leftarrow R_3-2R_2}
\begin{pmatrix}-1&3&2\\0&3&1\\0&0&0\end{pmatrix}
\end{gathered}
\]
\(\rank(A-I)=2\) (\lec{4}, Definition~10, Remark~3: the number of nonzero rows of an echelon form, first nonzero
entries not yet scaled to~\(1\)), \(g_A(1)=1\). \(-x_1+3x_2+2x_3=0\), \(3x_2+x_3=0\) (\lec{2},
Theorem~3(b)); \(x_3=3\): \(\alpha_1=(3,-1,3)'\), a basis of \(\mathcal N(A-I)\) (\lec{4}, Theorem~5(d)).
\end{board}
\begin{board}[\(\lambda=2\): reduce \(A-2I\); special solutions:]
\[
A-2I=\begin{pmatrix}3&-6&-6\\-1&2&2\\3&-6&-6\end{pmatrix}
\xrightarrow{R_1\leftrightarrow R_2}
\begin{pmatrix}-1&2&2\\3&-6&-6\\3&-6&-6\end{pmatrix}
\xrightarrow[R_3\leftarrow R_3+3R_1]{R_2\leftarrow R_2+3R_1}
\begin{pmatrix}-1&2&2\\0&0&0\\0&0&0\end{pmatrix}
\xrightarrow{R_1\leftarrow-R_1}
\begin{pmatrix}1&-2&-2\\0&0&0\\0&0&0\end{pmatrix}
\]
\(g_A(2)=3-1=2\); \(\alpha_2=(2,1,0)'\), \(\alpha_3=(2,0,1)'\), a basis of \(\mathcal N(A-2I)\) (\lec{4},
Theorem~8(c)).
\end{board}
\begin{board}[\(P\) and \(\Lambda\):]
\(g_A(1)+g_A(2)=3\): \(\alpha_1,\alpha_2,\alpha_3\) form a basis of \(\R^3\) (Theorem~\ref{h-thm-diag-criterion}(b)), so
\(P=[\alpha_1\ \alpha_2\ \alpha_3]\) is invertible (\lec{4}, Theorem~12(b),(d)), and \(AP=P\Lambda\) gives \(P^{-1}AP=\Lambda\)
(Theorem~\ref{h-thm-diag-criterion}(a)):
\[
P=\begin{pmatrix}3&2&2\\-1&1&0\\3&0&1\end{pmatrix},\qquad\Lambda=\diag(1,2,2),\qquad
AP=\begin{pmatrix}3&4&4\\-1&2&0\\3&0&2\end{pmatrix}=P\Lambda
\]
\end{board}

\Needspace{15\baselineskip}
\begin{theorem}\label{h-thm-powers}
If \(A=P\Lambda P^{-1}\) with \(P\in\R^{n\times n}\) invertible and \(\Lambda=\diag(\lambda_1,\dots,\lambda_n)\), then
\(A^k=P\Lambda^kP^{-1}=P\diag(\lambda_1^k,\dots,\lambda_n^k)P^{-1}\) for \(k=0,1,2,\dots\)
\end{theorem}
\begin{proof}
\begin{board}[Induction on \(k\); then \(\Lambda^k\):]
\boardtop
\begin{align*}
A^0&=I=P\Lambda^0P^{-1}\\
A^{k+1}&=A^kA=P\Lambda^kP^{-1}P\Lambda P^{-1}=P\Lambda^{k+1}P^{-1}&&\text{(\lec{1}, Theorem~3(a))}\\
\Lambda^k&=\diag(\lambda_1^k,\dots,\lambda_n^k)&&\text{(\lec{1}, Proposition~13(d))}
\end{align*}
\end{board}
\nopagebreak
\end{proof}

\Needspace{19\baselineskip}
\begin{example}\label{h-exm-diag-powers}
Find \(B^k\) (\(k=0,1,2,\dots\)) for the matrix \(B\) below. Multiplying out, \(P^{-1}P=I\) (so \(P\) is invertible,
\lec{2}, Theorem~15) and \(BP=P\Lambda\) with \(\Lambda=\diag(3,3,1,1)\):
\[
B=\begin{pmatrix}1&2&-2&2\\-2&3&-2&0\\0&-2&3&-2\\2&-2&2&1\end{pmatrix},\qquad
P=\begin{pmatrix}-1&0&0&-1\\0&-1&1&-1\\1&0&1&0\\0&1&0&1\end{pmatrix},\qquad
P^{-1}=\begin{pmatrix}0&-1&1&-1\\1&-1&1&0\\0&1&0&1\\-1&1&-1&1\end{pmatrix}
\]
\end{example}
\textit{Solution.} \(B=BPP^{-1}=P\Lambda P^{-1}\) (\lec{1}, Theorem~3(a)).
\begin{board}[\(B^k\) (Theorem~\ref{h-thm-powers}; \lec{1}, Theorem~3(a)), \(D=\diag(1,1,0,0)\), \(E=PDP^{-1}\):]
\boardtop
\begin{align*}
B^k&=P\diag(3^k,3^k,1,1)P^{-1}=P(3^kD+I-D)P^{-1}=3^kE+I-E&&(PP^{-1}=I)\\
B&=3E+I-E=I+2E,\quad\text{so }E=\tfrac12(B-I)&&(k=1)\\
B^k&=3^k\cdot\tfrac12(B-I)+I-\tfrac12(B-I)=I+\tfrac12(3^k-1)(B-I)&&\text{(\lec{1}, Definition~1)}\\
B^2&=I+4(B-I)=\begin{pmatrix}1&8&-8&8\\-8&9&-8&0\\0&-8&9&-8\\8&-8&8&1\end{pmatrix}
\end{align*}
\end{board}

\Needspace{19\baselineskip}
\subsection*{6. The products \(AB\) and \(BA\)}

\begin{theorem}\label{h-thm-AB-BA}
Let \(A\in\R^{m\times n}\) and \(B\in\R^{n\times m}\) with \(m\leqslant n\).
\begin{enumerate}
\item \(p_{BA}=t^{n-m}p_{AB}\): the eigenvalues of \(BA\), listed with multiplicity, are those of \(AB\) and \(n-m\)
zeros. So \(AB\), \(BA\) have the same nonzero eigenvalues, with the same algebraic multiplicities.
\item If \(c\ne0\) is a real eigenvalue of \(AB\) and the columns of \(X\in\R^{m\times g}\) form a basis of
\(\mathcal N(AB-cI_m)\), then the columns of \(BX\) are linearly independent eigenvectors of \(BA\) belonging to \(c\).
So \(g_{AB}(c)=g_{BA}(c)\).
\end{enumerate}
\end{theorem}
\begin{proof}
(a) For \(t\in\R\), \(t\ne0\), and \(M=\left(\begin{smallmatrix}I_n&t^{-1}B\\A&I_m\end{smallmatrix}\right)\), with
\(\det I=1\) (\lec{5}, Theorem~8(d)):
\begin{align*}
\det M&=\det I_n\det(I_m-AI_n^{-1}t^{-1}B)=\det(I_m-t^{-1}AB)&&\text{(\lec{5}, Theorem~19(a))}\\
\det M&=\det I_m\det(I_n-t^{-1}BI_m^{-1}A)=\det(I_n-t^{-1}BA)&&\text{(\lec{5}, Theorem~19(b))}\\
\det(tI_n-BA)&=t^n\det(I_n-t^{-1}BA)=t^n\det(I_m-t^{-1}AB)&&\text{(\lec{5}, Theorem~12(c))}\\
&=t^{n-m}\det(tI_m-AB)&&\text{(\lec{5}, Theorem~12(c))}
\end{align*}
So \(p_{BA}\) and \(t^{n-m}p_{AB}\) have the same value at every real \(t\ne0\) (\ref{f-values}), and \(p_{BA}=t^{n-m}p_{AB}\) (\ref{f-roots}). With \(p_{AB}=(t-\mu_1)\cdots(t-\mu_m)\)
(Definition~\ref{h-def-eigenvalues}), \(p_{BA}=(t-0)^{n-m}(t-\mu_1)\cdots(t-\mu_m)\): the eigenvalues of \(BA\) are
\(\mu_1,\dots,\mu_m\) and \(n-m\) zeros (\ref{f-complex-roots}), and \(a_{BA}(c)=a_{AB}(c)\) for \(c\ne0\)
(Definition~\ref{h-def-eigenvalues}).

(b) Let \(g=g_{AB}(c)\).
\begin{align*}
\rank X&=g,\qquad ABX=cX&&\text{(\lec{4}, Theorem~12(b); \lec{1}, Theorem~4(a))}\\
g=\rank X&\geqslant\rank(BX)\geqslant\rank(ABX)&&\text{(\lec{4}, Theorem~15(a))}\\
\rank(ABX)&=\rank\bigl(X(cI_g)\bigr)=\rank X=g&&\text{(\lec{4}, Theorem~15(d); \(cI_g\) invertible)}\\
BA(BX)&=B(ABX)=c\,BX&&\text{(\lec{1}, Theorem~3(a))}
\end{align*}
\begin{itemize}
\item So the \(g\) columns of \(BX\) are independent (\lec{4}, Theorem~12(b)), lie in \(\mathcal N(BA-cI_n)\)
(\lec{1}, Theorem~4(a)) and are nonzero (\lec{3}, Proposition~12(a)): eigenvectors of \(BA\) belonging to \(c\);
\(g_{BA}(c)\geqslant g\) (\lec{4}, Theorem~5(c): more than \(\dim W\) vectors of \(W\) are dependent).
\item The argument does not use \(m\leqslant n\); with \(A\), \(B\) exchanged it gives
\(g_{AB}(c)\geqslant g_{BA}(c)\).\qedhere
\end{itemize}
\end{proof}

\Needspace{17\baselineskip}
\begin{example}\label{h-exm-ata}
Find the eigenvalues of \(A'A\) and \(AA'\), and an eigenvector of \(AA'\) for each eigenvalue, where
\[
A=\begin{pmatrix}1&1&0\\1&0&1\\0&1&1\\-1&0&1\end{pmatrix},\qquad
A'A=\begin{pmatrix}3&1&0\\1&2&1\\0&1&3\end{pmatrix},\qquad
AA'=\begin{pmatrix}2&1&1&-1\\1&2&1&0\\1&1&2&1\\-1&0&1&2\end{pmatrix}
\]
\end{example}
\textit{Solution.}
\begin{board}[\(p_{A'A}\), expanding along row \(1\) (\lec{5}, Theorem~13(a)), and \(p_{AA'}\):]
\[
\begin{aligned}
\det(tI_3-A'A)&=\det\begin{pmatrix}t-3&-1&0\\-1&t-2&-1\\0&-1&t-3\end{pmatrix}
=(t-3)\bigl((t-2)(t-3)-1\bigr)-(t-3)\\
&=(t-3)(t^2-5t+4)=(t-1)(t-3)(t-4)
\end{aligned}
\]
Theorem~\ref{h-thm-AB-BA}(a), with \(A'\), \(A\) in place of \(A\), \(B\) (\(m=3\), \(n=4\)):
\(p_{AA'}=t\,p_{A'A}=t(t-1)(t-3)(t-4)\). Eigenvalues: \(1,3,4\) for \(A'A\); \(0,1,3,4\) for \(AA'\)
(Definition~\ref{h-def-eigenvalues}).
\end{board}
\Needspace{8\baselineskip}
\begin{board}[\(\lambda=4,3,1\): eigenvectors of \(A'A\), mapped by \(A\) (Theorem~\ref{h-thm-AB-BA}(b) with \(A'\),
\(A\) in place of \(A\), \(B\)):]
\boardtop
\begin{align*}
A'A(1,1,1)'&=(4,4,4)'=4(1,1,1)',&A(1,1,1)'&=(2,2,2,0)'\\
A'A(-1,0,1)'&=(-3,0,3)'=3(-1,0,1)',&A(-1,0,1)'&=(-1,0,1,2)'\\
A'A(1,-2,1)'&=(1,-2,1)'=1(1,-2,1)',&A(1,-2,1)'&=(-1,2,-1,0)'
\end{align*}
\(4\), \(3\), \(1\) are simple, so each of the three vectors \(v\) is a basis of \(\mathcal N(A'A-\lambda I_3)\)
(Theorem~\ref{h-thm-geo-alg}; \lec{4}, Theorem~5(d)); with \(X=v\): \((2,2,2,0)'\), \((-1,0,1,2)'\), \((-1,2,-1,0)'\) are
eigenvectors of \(AA'\) for \(4\), \(3\), \(1\).
\end{board}
\begin{board}[\(\lambda=0\): \(\mathcal N(AA')=\mathcal N(A')\) (\lec{4}, Theorem~16(a) with \(A'\)); reduce \(A'\):]
\[
A'=\begin{pmatrix}1&1&0&-1\\1&0&1&0\\0&1&1&1\end{pmatrix}
\xrightarrow{R_2\leftarrow R_2-R_1}
\begin{pmatrix}1&1&0&-1\\0&-1&1&1\\0&1&1&1\end{pmatrix}
\xrightarrow{R_3\leftarrow R_3+R_2}
\begin{pmatrix}1&1&0&-1\\0&-1&1&1\\0&0&2&2\end{pmatrix}
\]
\(\rank A'=3\), \(\dim\mathcal N(A')=4-3=1\) (\lec{4}, Definition~10, Remark~3, Theorem~11). Back substitution (\lec{2},
Theorem~3(b)), \(x_4=1\): \(x_3=-1\), \(x_2=0\), \(x_1=1\); \((1,0,-1,1)'\) is a basis of \(\mathcal N(AA')\) (\lec{4},
Theorem~5(d)).
\end{board}

\bigskip
\textbf{Practice.} The Exercises and Problems of Chapter~7 of the textbook.

\end{document}
